% ============================================================================
%  Atlas - SCRIPT DE DEMONSTRATION du POC (a derouler devant un jury)
%  Compilation : pdflatex POC_Atlas_Demo.tex  (2 passes pour la TOC)
%  But : la sequence EXACTE de commandes, dans l'ordre, pour MONTRER le POC.
%  Chaque acte = une commande a lancer + ce que le public doit voir.
% ============================================================================
\documentclass[11pt,a4paper]{article}

\usepackage[T1]{fontenc}
\usepackage[utf8]{inputenc}
\usepackage{lmodern}
\usepackage[french]{babel}
\usepackage[a4paper,margin=2.3cm]{geometry}
\usepackage{xcolor}
\usepackage{booktabs}
\usepackage{tabularx}
\usepackage{array}
\usepackage{listings}
\usepackage[most]{tcolorbox}
\usepackage{titlesec}
\usepackage{fancyhdr}
\usepackage{enumitem}
\usepackage{underscore}
\usepackage[hidelinks,colorlinks=true,linkcolor=atlasblue,urlcolor=atlasblue]{hyperref}

\definecolor{atlasblue}{HTML}{1F6FEB}
\definecolor{atlasdark}{HTML}{0D1B2A}
\definecolor{atlasgray}{HTML}{F2F4F8}
\definecolor{atlasgreen}{HTML}{2EA043}
\definecolor{atlasred}{HTML}{CF222E}
\definecolor{atlasamber}{HTML}{9A6700}
\definecolor{codebg}{HTML}{F6F8FA}
\definecolor{codecom}{HTML}{6E7781}
\definecolor{codestr}{HTML}{0A3069}

\titleformat{\section}{\Large\bfseries\color{atlasblue}}{}{0em}{}
\titlespacing*{\section}{0pt}{2.4ex plus 1ex minus .2ex}{1.2ex}

\pagestyle{fancy}
\fancyhf{}
\lhead{\small\color{atlasdark}\textbf{Atlas} — Script de démonstration}
\rhead{\small\color{codecom}\thepage}
\rfoot{\small\color{codecom}Support de présentation — POC}
\renewcommand{\headrulewidth}{0.4pt}

\lstset{
  basicstyle=\ttfamily\small,
  backgroundcolor=\color{codebg},
  frame=leftline, framerule=2.5pt, rulecolor=\color{atlasblue},
  keywordstyle=\color{atlasblue}\bfseries,
  commentstyle=\color{codecom}\itshape,
  stringstyle=\color{codestr},
  showstringspaces=false, breaklines=true, breakatwhitespace=true,
  columns=fullflexible, keepspaces=true, xleftmargin=1em,
  aboveskip=0.8em, belowskip=0.8em, inputencoding=utf8, extendedchars=true,
  literate=%
    {é}{{\'e}}1 {è}{{\`e}}1 {ê}{{\^e}}1 {ë}{{\"e}}1
    {à}{{\`a}}1 {â}{{\^a}}1 {ä}{{\"a}}1
    {î}{{\^i}}1 {ï}{{\"i}}1 {ô}{{\^o}}1 {ö}{{\"o}}1
    {û}{{\^u}}1 {ù}{{\`u}}1 {ü}{{\"u}}1 {ç}{{\c c}}1
    {É}{{\'E}}1 {È}{{\`E}}1 {À}{{\`A}}1
    {«}{{\guillemotleft}}1 {»}{{\guillemotright}}1
    {→}{{$\rightarrow$}}1 {≤}{{$\leq$}}1
}
\lstdefinestyle{bash}{language=bash,morekeywords={ansible,ansible-playbook,ssh,pct,mysql}}

% Boite "ce qu'on montre / ce que le jury voit"
\newtcolorbox{montrer}{enhanced,breakable,colback=atlasgreen!7,colframe=atlasgreen,
  fonttitle=\bfseries,title=\faEyeText\ À l'écran — ce que le public voit,
  left=2mm,right=2mm,top=1mm,bottom=1mm}
% Boite "phrase a dire"
\newtcolorbox{dire}{enhanced,breakable,colback=atlasblue!6,colframe=atlasblue,
  fonttitle=\bfseries,title=Ce qu'on dit,left=2mm,right=2mm,top=1mm,bottom=1mm}
\newtcolorbox{attention}{enhanced,breakable,colback=atlasred!6,colframe=atlasred,
  fonttitle=\bfseries,title=Attention,left=2mm,right=2mm,top=1mm,bottom=1mm}
% Bandeau d'acte : \acte{titre}{duree}
\newcommand{\acte}[2]{%
  \begin{tcolorbox}[enhanced,colback=atlasdark,colframe=atlasdark,coltext=white,
    sharp corners,boxrule=0pt,left=3mm,right=3mm,top=2mm,bottom=2mm,
    before skip=10pt,after skip=6pt]%
    {\bfseries\large #1}\dur{#2}%
  \end{tcolorbox}}

\newcommand{\faEyeText}{\raisebox{-1pt}{$\odot$}}
\newcommand{\cmd}[1]{{\ttfamily #1}}
\newcommand{\dur}[1]{\hfill{\small\color{codecom}\itshape (#1)}}

\begin{document}

% ===================== TITRE =====================
\begin{titlepage}
\centering
\vspace*{2cm}
{\color{atlasblue}\rule{\linewidth}{2pt}}\\[0.6cm]
{\Huge\bfseries\color{atlasdark} Projet Atlas}\\[0.3cm]
{\LARGE Script de démonstration du POC}\\[0.2cm]
{\large\color{codecom} La séquence exacte de commandes, dans l'ordre, pour présenter}\\[0.4cm]
{\color{atlasblue}\rule{\linewidth}{2pt}}\\[1.2cm]

\begin{tcolorbox}[width=0.9\linewidth,colback=atlasgray,colframe=atlasdark,sharp corners]
\centering\small
\textbf{Fil rouge de la démo} : \\[3pt]
\textbf{1.} Ça tourne \; $\rightarrow$ \; \textbf{2.} C'est du code (IaC) \; $\rightarrow$ \;
\textbf{3.} C'est supervisé \& résilient \\[2pt]
$\rightarrow$ \; \textbf{4.} Les données sont protégées (RPO) \; $\rightarrow$ \;
\textbf{5.} On survit au désastre (RTO)\\[6pt]
\textbf{Durée cible} : 12–15 min \quad|\quad \textbf{Objectifs prouvés} : IaC, sécurité, supervision, RPO $\leq$ 20 min, RTO $\leq$ 40 min
\end{tcolorbox}

\vfill
{\color{codecom}\small Support de présentation — à garder à côté de soi pendant la démo}\\
{\color{codecom}\small Généré le \today}
\end{titlepage}

\tableofcontents
\newpage

% ===================== PREPARATION =====================
\section*{Avant de commencer (préparation, hors public)}
\addcontentsline{toc}{section}{Avant de commencer}
Ouvrir \textbf{à l'avance} 4 onglets/terminaux pour enchaîner sans temps mort :

\begin{lstlisting}[style=bash]
# Terminal A : sur le control node (hyperviseur), pour Ansible + pct
ssh root@138.201.135.108

# Terminaux B, C, D : les 3 tunnels (a laisser ouverts), sur son poste
ssh -L 8443:10.20.0.10:443  root@138.201.135.108   # GLPI
ssh -L 8080:10.30.0.12:8080 root@138.201.135.108   # Zabbix
ssh -L 3000:10.30.0.14:3000 root@138.201.135.108   # Grafana
\end{lstlisting}
Puis ouvrir 3 onglets navigateur : \url{https://localhost:8443}, \url{http://localhost:8080},
\url{http://localhost:3000}. \textbf{Ne pas encore les montrer.}

\subsection*{« Armer » l'incident chaos (indispensable)}
Le trigger disque de Zabbix exige que le disque reste $>$90\% \textbf{pendant 5 minutes}
(anti-faux-positif) : l'alerte n'apparaît donc que \textbf{5 min après} le déclenchement.
On \textbf{arme donc l'incident tout de suite}, avant l'Acte 1 ; il sera visible pile au
moment de l'Acte 3.

\begin{lstlisting}[style=bash]
# (Terminal A, TOUT AU DEBUT) declencher le remplissage disque sur la machine de lab
ssh root@138.201.135.108 'pct exec 208 -- /opt/atlas/scripts/chaos.sh disk'
\end{lstlisting}

\begin{attention}
Vérifier 5 min avant : les 3 tunnels montent, Zabbix affiche 10 hôtes verts (état sain de
départ), et la machine de lab est propre (\cmd{... chaos.sh clean} si un incident cron
traîne). \textbf{Puis} armer l'incident disque ci-dessus juste avant de commencer à parler.
\end{attention}

% ===================== ACTE 1 =====================
\acte{ACTE 1 — « Ça tourne » : le service est vivant}{~3 min}

\subsection*{1.1 — Toute l'infra est debout}
\begin{lstlisting}[style=bash]
# (Terminal A, sur le control node)
ssh root@138.201.135.108 'pct list'
\end{lstlisting}
\begin{montrer}
Les \textbf{9 conteneurs LXC} sont tous en statut \cmd{running} (bastion, reverse proxy,
GLPI, MariaDB, SMTP, Zabbix, backup, Grafana, testlab).
\end{montrer}
\begin{dire}
« Voici les 9 briques du POC, réparties sur 3 VLAN — DMZ, LAN-admin, management. Tout
tourne. Regardons le service côté utilisateur. »
\end{dire}

\subsection*{1.2 — Le portail GLPI répond (chemin utilisateur, en HTTPS)}
\begin{lstlisting}[style=bash]
# Onglet navigateur GLPI (tunnel deja ouvert)
https://localhost:8443
\end{lstlisting}
\begin{montrer}
La page GLPI s'affiche en \textbf{HTTPS} (via le reverse proxy TLS). Se connecter, montrer
la \textbf{liste des tickets} existants.
\end{montrer}
\begin{dire}
« Le trafic utilisateur entre par le reverse proxy chiffré, qui sert GLPI — la base
n'est jamais exposée directement. »
\end{dire}

\subsection*{1.3 — Tout est supervisé (Zabbix)}
\begin{lstlisting}[style=bash]
http://localhost:8080     # Monitoring -> Hosts
\end{lstlisting}
\begin{montrer}
\textbf{10 hôtes en vert} (\cmd{Available}) et \emph{Monitoring $\rightarrow$ Problems}
\textbf{vide}. C'est notre état de référence « tout va bien ».
\end{montrer}

\subsection*{1.4 — Les tableaux de bord (Grafana)}
\begin{lstlisting}[style=bash]
http://localhost:3000     # dashboard "Atlas - Vue d'ensemble"
\end{lstlisting}
\begin{montrer}
Grafana remonte les métriques (via un compte Zabbix en \textbf{lecture seule}). On voit
CPU/RAM/disque des hôtes en direct.
\end{montrer}

% ===================== ACTE 2 =====================
\acte{ACTE 2 — « C'est du code » : Infrastructure as Code}{~2 min}

\subsection*{2.1 — Rien n'est fait à la main : tout est reproductible}
\begin{lstlisting}[style=bash]
# (Terminal A) simulation : est-ce que le reel colle au code ?
ansible-playbook site.yaml --check --diff --tags bastion --limit atlbst01p
\end{lstlisting}
\begin{montrer}
Le playbook tourne en \textbf{mode simulation} et le bilan affiche
\cmd{changed=0  failed=0} : l'infrastructure réelle est \textbf{déjà conforme au code}.
Aucune dérive.
\end{montrer}
\begin{attention}
Utiliser le tag \cmd{bastion} (et non \cmd{zabbix}) : en mode \cmd{--check}, le rôle
\cmd{zabbix} échoue faussement (une tâche crée le dossier TLS, simulé donc absent, et la
tâche suivante ne le trouve pas). Le rôle \cmd{bastion} donne un \cmd{changed=0} propre.
\end{attention}
\begin{dire}
« Toute l'infra est décrite en Ansible. Je peux la rejouer à l'identique : ici la
simulation ne propose aucun changement, preuve que le serveur EST son code. C'est ce qui
garantit qu'on saura tout reconstruire à l'identique en cas de sinistre. »
\end{dire}

% ===================== ACTE 3 =====================
\acte{ACTE 3 — « C'est résilient » : chaos en direct}{~3 min}

\subsection*{3.1 — Rappel : l'incident a été armé au début}
On a lancé le remplissage disque en préparation (\cmd{chaos.sh disk} sur \cmd{atltst01p}) :
le disque est monté à $>$90\%. \textbf{5 minutes plus tard}, le trigger se déclenche — c'est
maintenant.
\begin{dire}
« En début de démo, j'ai provoqué une panne réelle mais contrôlée sur la machine de lab :
son disque s'est rempli. Voyons ce que la supervision en a fait, sans que j'aie touché à
Zabbix. »
\end{dire}

\subsection*{3.2 — Zabbix a détecté la panne tout seul}
\begin{lstlisting}[style=bash]
http://localhost:8080     # Monitoring -> Problems (rafraichir)
\end{lstlisting}
\begin{montrer}
Un problème \textbf{rouge (Haute sévérité)} sur \cmd{atltst01p} :
\textbf{« FS [/]: Space is critically low (used > 90\%) »}. Détection automatique, aucune
config en direct. \emph{(Validé en répétition : alerte levée $\approx$ 5 min après le
déclenchement.)}
\end{montrer}

\subsection*{3.3 — On répare, et l'alerte se résout}
\begin{lstlisting}[style=bash]
# (Terminal A)
ssh root@138.201.135.108 'pct exec 208 -- /opt/atlas/scripts/chaos.sh clean'
\end{lstlisting}
\begin{montrer}
Le disque de \cmd{atltst01p} revient à $\sim$24\%. Dans Zabbix, le problème repasse en
\textbf{Resolved} et l'hôte redevient vert : la boucle « je casse $\rightarrow$ ça alerte
$\rightarrow$ je répare » est bouclée.
\end{montrer}
\begin{attention}
\cmd{chaos.sh} a un garde-fou : il refuse de s'exécuter ailleurs que sur \cmd{atltst01p}
(vérif \cmd{hostname -s}). Aucun risque pour les vrais services. \textbf{Ne pas utiliser le
scénario \cmd{service}} (arrêt de nginx) : nginx n'est pas supervisé sur la machine de lab,
donc il ne lèverait aucune alerte. Le scénario \cmd{disk} est le bon pour la démo.
\end{attention}

% ===================== ACTE 4 =====================
\acte{ACTE 4 — « Les données sont protégées » : le RPO}{~2 min}

\subsection*{4.1 — Sauvegarde toutes les 15 minutes}
\begin{lstlisting}[style=bash]
# (Terminal A) les derniers dumps sur la base
ssh root@10.30.0.10 'ls -lht /var/backups/atlas/db/ | head'
\end{lstlisting}
\begin{montrer}
Une série de dumps espacés de \textbf{15 minutes} (\cmd{glpi_dump_AAAAMMJJ_HHMMSS.sql.gz}),
le plus récent daté de moins de 15 min $\rightarrow$ \textbf{RPO $\leq$ 15 min}
(mieux que la cible de 20 min).
\end{montrer}

\subsection*{4.2 — Répliqué hors-site}
\begin{lstlisting}[style=bash]
ssh root@10.30.0.13 'ls -lht /var/backups/atlas/db-remote/ | head'
\end{lstlisting}
\begin{montrer}
Le même dernier dump est présent sur le serveur de sauvegarde \textbf{distinct}
(\cmd{atlbkp01p}) : la perte de la base n'emporte pas les sauvegardes.
\end{montrer}

% ===================== ACTE 5 =====================
\acte{ACTE 5 — « On survit au désastre » : le PRA (RTO)}{~3–4 min}

\subsection*{5.1 — Option courte à montrer en direct : restauration granulaire}
\begin{lstlisting}[style=bash]
# (Dans GLPI) supprimer un ticket de demonstration, puis :
# (Terminal A) choisir un dump anterieur
ssh root@10.30.0.10 'ls -lht /var/backups/atlas/db/ | head'

# Relancer uniquement les tickets manquants (sans ecraser les autres)
ansible-playbook pra_restore_granular.yaml \
  -e "pra_target_dump=/var/backups/atlas/db/glpi_dump_AAAAMMJJ_HHMMSS.sql.gz"
\end{lstlisting}
\begin{montrer}
Le playbook restaure dans une base temporaire, réinjecte \emph{seulement} les tickets
absents (\cmd{INSERT IGNORE}), affiche le nombre de tickets, et journalise le RPO effectif.
Rafraîchir GLPI : \textbf{le ticket supprimé est revenu}, les autres sont intacts.
\end{montrer}
\begin{dire}
« On peut réparer chirurgicalement une erreur humaine, sans toucher au reste. »
\end{dire}

\subsection*{5.2 — Option complète (à décrire, ou à lancer si le temps le permet) : perte de VM}
\begin{lstlisting}[style=bash]
ansible-playbook pra_restore_full.yaml
\end{lstlisting}
\begin{montrer}
Reconstruction de toute la chaîne dans l'ordre (MariaDB $\rightarrow$ GLPI $\rightarrow$
reverse proxy $\rightarrow$ supervision), restauration des données, vérifications
automatiques (GLPI en HTTPS, comptage des tickets), et surtout : à la fin, le playbook
\textbf{affiche le RTO mesuré} comparé à la cible de 40 min.
\end{montrer}
\begin{dire}
« En une seule commande, tout se reconstruit et se vérifie tout seul, et le temps de reprise
est mesuré : c'est ça, tenir le RTO. »
\end{dire}

% ===================== CLOTURE =====================
\acte{Clôture — ce qu'on a prouvé}{~1 min}

\begin{table}[h]
\centering\small
\renewcommand{\arraystretch}{1.25}
\begin{tabularx}{\linewidth}{@{}l X@{}}
\toprule
\textbf{Exigence} & \textbf{Démontré par} \\
\midrule
100\% IaC, reproductible   & Acte 2 — simulation « 0 changed », infra = code \\
Sécurité (segmentation)    & Acte 1 — HTTPS via reverse proxy, base non exposée, 3 VLAN \\
Supervision opérationnelle & Acte 3 — panne détectée puis résolue en direct \\
RPO $\leq$ 20 min          & Acte 4 — dumps toutes les 15 min, répliqués hors-site \\
RTO $\leq$ 40 min          & Acte 5 — restauration en une commande, RTO mesuré \\
\bottomrule
\end{tabularx}
\end{table}

\begin{center}
{\color{atlasblue}\rule{0.5\linewidth}{1pt}}\\[4pt]
{\small\color{codecom} Atlas — Script de démonstration du POC — Infrastructure as Code}
\end{center}

% ===================== AIDE-MEMOIRE 1 PAGE =====================
\newpage
\section*{Aide-mémoire — la démo en une page (à imprimer)}
\addcontentsline{toc}{section}{Aide-mémoire (1 page)}

\begin{enumerate}[leftmargin=1.5em,itemsep=4pt]
  \item \textbf{Au tout début — armer le chaos} : \cmd{ssh root@138.201.135.108 'pct exec 208 -- /opt/atlas/scripts/chaos.sh disk'} (l'alerte arrive $\approx$5 min après, pile pour l'Acte 3).
  \item \textbf{Ça tourne} : \cmd{ssh root@138.201.135.108 'pct list'} $\rightarrow$ 9 running ;
        puis les 3 onglets : GLPI (HTTPS), Zabbix (10 verts), Grafana.
  \item \textbf{C'est du code} : \cmd{ansible-playbook site.yaml --check --diff --tags bastion --limit atlbst01p} $\rightarrow$ changed=0.
  \item \textbf{Chaos} : Zabbix $\rightarrow$ Problems affiche « FS [/]: Space is critically low » (rouge) $\rightarrow$ réparer : \cmd{ssh ... 'pct exec 208 -- .../chaos.sh clean'} $\rightarrow$ résolu.
  \item \textbf{RPO} : \cmd{ssh root@10.30.0.10 'ls -lht /var/backups/atlas/db/ | head'} + \cmd{ssh root@10.30.0.13 '... db-remote ...'}.
  \item \textbf{RTO} : \cmd{ansible-playbook pra_restore_granular.yaml -e "pra_target_dump=..."} (ou \cmd{pra_restore_full.yaml}).
\end{enumerate}

\end{document}
